\documentclass[12pt,a4paper]{book}
%\documentclass[a4paper,french]{book}
%\usepackage[17pt]{extsizes}


\input{../1ere-preambule}

\rhead{DS n$^{\circ}04$ - Dérivation - 14/04/2025}

\newcommand{\va}[1]{\lvert \ #1 \ \rvert}
\newcommand{\vaf}[1]{\left\lvert \ #1\  \right\rvert}

\begin{document}
	\graphicspath{{images/}}
	
	\begin{center}
		
		
		\fbox{
			\begin{minipage}[]{15cm}
				\begin{center}
					\LARGE{\rule{0.5cm}{0pt}\rule[-0.25cm]{0pt}{0.9cm}
						DS n$^{\circ}04$ - 14/04/2025 \\
						Dérivation 
					}
				\end{center}
			\end{minipage}
		}
	\end{center}
	
	\begin{center}\textbf{1h20min - Barème indicatif}\end{center}
	\begin{center}\textbf{Les élèves avec un tier-temps ne traitent pas les questions avec  le symbole \ding{46}}\end{center}
	\begin{center}\textbf{Calculatrice autorisée.\\Les résultats devront être justifiés (à l'aide de calcul ou de propriétés).}\end{center}

\exerciceDS{ - (1,5 points)}

\begin{minipage}[c]{0.48\linewidth}
	A l'aide du graphique ci-contre et sachant que la courbe verte est la courbe représentative de la fonction f et que les droites rouges sont les tangentes à la courbe en 2 ; 4 et 8; déterminer $f'(2)$, $f'(4)$ et $f'(8)$.
\end{minipage} \hfill
\begin{minipage}[c]{0.48\linewidth}
	\begin{center}
		\newrgbcolor{qqwuqq}{0. 0.39215686274509803 0.}
		\psset{xunit=0.8cm,yunit=0.8cm,algebraic=true,dimen=middle,dotstyle=o,dotsize=5pt 0,linewidth=1.5pt,arrowsize=3pt 2,arrowinset=0.25}
		\begin{pspicture*}(-0.7,-9.2)(9.,0.5)
%			\multips(0,-5)(0,1.0){11}{\psline[linestyle=dashed,linecap=1,dash=1.5pt 1.5pt,linewidth=0.4pt,linecolor=lightgray]{c-c}(-2.5,0)(7.,0)}
%			\multips(-2,0)(1.0,0){10}{\psline[linestyle=dashed,linecap=1,dash=1.5pt 1.5pt,linewidth=0.4pt,linecolor=lightgray]{c-c}(0,-5.)(0,5.2)}
			\psgrid[subgriddiv=0,gridcolor=gray, gridlabels=0pt](-2,-9)(9,1)
			\psaxes[labelFontSize=\scriptstyle,xAxis=true,yAxis=true,Dx=1.,Dy=1.,ticksize=-2pt 0,subticks=2]{->}(0,0)(-2.5,-9.)(9.,0.5)
			\psplot[linewidth=1.2pt,linecolor=qqwuqq,plotpoints=200]{-2.5}{9.0}{0.25*x^(2.0)-2.0*x-2.0}
			\psplot[linewidth=1.2pt,linecolor=red]{-2.5}{9.}{(-3.-1.*x)/1.}
			\psplot[linewidth=1.2pt,linecolor=red]{-2.5}{9.}{(-6.-0.*x)/1.}
			\psplot[linewidth=1.2pt,linecolor=red]{-2.5}{9.}{(-18.--2.*x)/1.}
			%		\begin{scriptsize}
				\rput[bl](-3.52,-11.04){\qqwuqq{$f$}}
				\rput[bl](2.22,5.82){\red{$g$}}
				\rput[bl](-4.14,3.68){\red{$h$}}
				\rput[bl](3.62,5.82){\red{$i$}}
				%		\end{scriptsize}
		\end{pspicture*}
	\end{center}
\end{minipage}





\exerciceDS{ - (2 points)}

\begin{minipage}[c]{0.48\linewidth}
	Dans le graphique ci-contre, tracer les tangentes en 1 et en 3 à la courbe représentative de la fonction f sachant que $f'(1)=2$ et $f'(3)=-6$
\end{minipage} \hfill
\begin{minipage}[c]{0.48\linewidth}
	\begin{center}
		\psset{xunit=0.8cm,yunit=0.8cm,algebraic=true,dimen=middle,dotstyle=o,dotsize=5pt 0,linewidth=1.2pt,arrowsize=3pt 2,arrowinset=0.25}
		\begin{pspicture*}(-2.,-2)(5.,6)
			\psgrid[subgriddiv=0,gridcolor=gray, gridlabels=0pt](-2.,-2)(5.,6)
%			\multips(0,-3)(0,1.0){9}{\psline[linestyle=dashed,linecap=1,dash=1.5pt 1.5pt,linewidth=0.4pt,linecolor=lightgray]{c-c}(-2.,0)(5.,0)}
%			\multips(-2,0)(1.0,0){8}{\psline[linestyle=dashed,linecap=1,dash=1.5pt 1.5pt,linewidth=0.4pt,linecolor=lightgray]{c-c}(0,-3.5)(0,4.5)}
			\psaxes[labelFontSize=\scriptstyle,xAxis=true,yAxis=true,Dx=1.,Dy=1.,ticksize=-2pt 0,subticks=2]{->}(0,0)(-2.,-2)(5.,6)
			\psplot[linewidth=1.2pt,linecolor=ForestGreen,plotpoints=200]{-2.0}{5.0}{-2.0*x^(2.0)+6.0*x}
		\end{pspicture*}
	\end{center}
\end{minipage}

\newpage

\exerciceDS{ - (4 points)}

Déterminer sur $\mathbb{R}$ l'expression de la fonction dérivée de chacune des fonctions suivantes.
\begin{enumMet}
	\item $f(x)= 3x^2-7 x+4$
	\item $g(x)= 7 x^3+4 x^2-7$
	\item $h(x)= -4x^3+6x^2-3x+2$
	\item $k(x)=-6 x^3+7 x-5x^2+3$
\end{enumMet}

\exerciceDS{ - (4 points)}

Déterminer sur $\mathbb{R}$ l'expression de la fonction dérivée de chacune des fonctions suivantes.
\begin{enumMet}
	\item $f(x)= 0,9 x^2-7,2 x+4$
	\item $g(x)= 3-6 x^2$
	\item $h(x)= 5 x^3-13 x-15$
	\item $k(x)=-\dfrac{5}{3} x^3-\dfrac{7}{2} x^2$
\end{enumMet}

\exerciceDS{ - (4 points)}

Soit $f$ la fonction définie sur \R par $f(x)=-2x^2+5x-4$
\begin{enumMet}
	\item Déterminer $f'$ la fonction dérivée de $f$ sur \R
	\item Etudier le signe de $f'(x)$ sur \R et en déduire les variations de la fonction $f$ sur \R
\end{enumMet}

\exerciceDS{ - (5 points)}

Soit $f$ la fonction définie sur \R par $f(x)=-2x^3+\dfrac{19}{2}x^2-10x+2$
\begin{enumMet}
	\item Déterminer $f'$ la fonction dérivée de $f$ sur \R
	\item Montrer que $f'(x)=(-3x+2)(2x-5)$
	\item Etudier le signe de $f'(x)$ sur \R et en déduire les variations de la fonction $f$ sur \R
\end{enumMet}





\end{document}

